\section*{Chapter \ref{ch:b1:crystals-ionic-covalent} --- Crystals II: Ionic, Covalent and Molecular Solids}

\begin{solution}{exo:b1:crystals-ionic-covalent:1}
Anions: $8 \times \frac18 + 6 \times \frac12 = 4$; cations: $12 \times
\frac14 + 1 = 4$. Each ion has 6 neighbours of the other kind (6:6). Four
cations and four anions: \ce{NaCl}.
\end{solution}

\begin{solution}{exo:b1:crystals-ionic-covalent:2}
Half the main diagonal: $a\sqrt3/2 = 412.3 \times 0.866 = \qty{357.1}{pm}$.
One \ce{Cs+} and $8 \times \frac18 = 1$ \ce{Cl-}: one formula unit.
\end{solution}

\begin{solution}{exo:b1:crystals-ionic-covalent:3}
$x = 102/181 = 0.56$, between 0.414 and 0.732: octahedral coordination
(6), the \omterm{def:b1:crystals-ionic-covalent:structure-type}{rock-salt type}.
\end{solution}

\begin{solution}{exo:b1:crystals-ionic-covalent:4}
Metallic: copper. Ionic: sodium chloride, calcium fluoride. Covalent:
diamond, quartz. Molecular: diiodine, ice.
\end{solution}

\begin{solution}{exo:b1:crystals-ionic-covalent:5}
Along an edge, anion--cation--anion in contact: $a = 2(r_+ + r_-)$. Along
a face diagonal lie two anions at the distance $a\sqrt2/2$; they must not
overlap: $a\sqrt2/2 \ge 2r_-$, so $(r_+ + r_-)\sqrt2 \ge 2r_-$, that is
$r_+/r_- \ge \sqrt2 - 1$.
\end{solution}

\begin{solution}{exo:b1:crystals-ionic-covalent:6}
$\rho = 4 \times 58.44/(6.022 \times 10^{23} \times (5.641 \times
10^{-8})^3) = \qty{2.16}{g/cm^3}$, against \qty{2.17}{g/cm^3} measured.
\end{solution}

\begin{solution}{exo:b1:crystals-ionic-covalent:7}
$\rho = 168.36/(6.022 \times 10^{23} \times (4.123 \times 10^{-8})^3) =
\qty{3.99}{g/cm^3}$, the measured value.
\end{solution}

\begin{solution}{exo:b1:crystals-ionic-covalent:8}
\ce{Zn-S} $= a\sqrt3/4 = \qty{234.2}{pm}$;
$\rho = 4 \times 97.44/(6.022 \times 10^{23} \times (5.409 \times
10^{-8})^3) = \qty{4.09}{g/cm^3}$, close to the measured
\qty{4.04}{g/cm^3} of natural sphalerite.
\end{solution}

\begin{solution}{exo:b1:crystals-ionic-covalent:9}
\ce{C-C} $= 245.6/\sqrt3 = \qty{141.8}{pm}$; interlayer $669.6/2 =
\qty{334.8}{pm}$. Cell volume $\frac{\sqrt3}{2}a^2c = \qty{3.498e-23}{cm^3}$,
$\rho = 4 \times 12.01/(6.022 \times 10^{23} \times 3.498 \times 10^{-23})
= \qty{2.28}{g/cm^3}$, against \qty{3.52}{g/cm^3} for diamond: the layers
are far apart, held only by London forces.
\end{solution}

\begin{solution}{exo:b1:crystals-ionic-covalent:10}
See the proof of \cref{prop:b1:crystals-ionic-covalent:diamond}:
$C = \pi\sqrt3/16 = 0.34$. Hardness comes from the strength and the
three-dimensional rigidity of the \omterm{def:b1:lewis-resonance-vsepr:lewis}{covalent bonds}, not from packing: each
atom has only four neighbours, at fixed angles, which leaves much empty
space.
\end{solution}

\begin{solution}{exo:b1:crystals-ionic-covalent:11}
\ce{Si-Si} $= a\sqrt3/4 = \qty{235.2}{pm}$, \omterm{def:b1:periodicity:radius}{covalent radius}
\qty{117.6}{pm}; $\rho = 8 \times 28.09/(6.022 \times 10^{23} \times
(5.431 \times 10^{-8})^3) = \qty{2.33}{g/cm^3}$, the measured value.
\end{solution}

\begin{solution}{exo:b1:crystals-ionic-covalent:12}
$x = 152/181 = 0.84 > 0.732$: the rule predicts the \omterm{def:b1:crystals-ionic-covalent:structure-type}{caesium chloride type}.
In the rock-salt structure, \ce{Rb-Cl} $= a/2 = \qty{329.1}{pm}$, close to
$152 + 181 = \qty{333}{pm}$: the observed structure is consistent. The
rule treats ions as hard spheres and ignores their \omterm{def:b1:periodicity:polarisability}{polarisability} and
the small difference of energy between the two structures (rubidium
chloride takes the \omterm{def:b1:crystals-ionic-covalent:structure-type}{caesium chloride type} under pressure); it is a guide,
not a law.
\end{solution}

\begin{solution}{pb:b1:crystals-ionic-covalent:1}
\textbf{1.} At the 8 corners and the 6 face centres: $8 \times \frac18 + 6
\times \frac12 = 4$ ions.
\textbf{2.} 8 \omterm{def:b1:crystals-metals:site}{tetrahedral sites}, all occupied: 8 \ce{F-} for 4
\ce{Ca^{2+}}, ratio 1:2, \ce{CaF2}.
\textbf{3.} \ce{F-}: 4 \ce{Ca^{2+}} at the corners of a tetrahedron.
\ce{Ca^{2+}}: 8 \ce{F-} at the corners of a cube.
\textbf{4.} $a\sqrt3/4 = 546.3 \times 0.4330 = \qty{236.6}{pm}$.
\textbf{5.} $112 + 131 = \qty{243}{pm}$: within 3\,\%.
\textbf{6.} The \ce{F-} ions form a simple cubic array of edge $a/2 =
\qty{273.2}{pm}$, slightly more than $2 \times 131 = \qty{262}{pm}$: they
almost touch.
\textbf{7.} $x = 112/131 = 0.85$.
\textbf{8.} $x > 0.732$: coordination 8.
\textbf{9.} The \omterm{def:b1:crystals-ionic-covalent:structure-type}{caesium chloride type} has as many cations as anions; with
twice as many anions, only half of the cubes of anions can hold a cation.
\textbf{10.} The 8 \ce{F-} of the cell form a cube of 8 small cubes of edge
$a/2$; the \ce{Ca^{2+}} ions occupy the centres of 4 of them, in
alternation, like the black squares of a three-dimensional chessboard.
\textbf{11.} $C = \frac43\pi(4 \times 112^3 + 8 \times 131^3)/546.3^3 = 0.61$.
\textbf{12.} \ce{O^{2-}} on the FCC \omterm{def:b1:crystals-metals:lattice}{lattice} (corners and face centres),
\ce{Li+} in the 8 \omterm{def:b1:crystals-metals:site}{tetrahedral sites}.
\textbf{13.} $a\sqrt3/4 = \qty{200.0}{pm}$; $59 + 142 = \qty{201}{pm}$.
\textbf{14.} $\rho = 4 \times 29.88/(6.022 \times 10^{23} \times (4.619
\times 10^{-8})^3) = \qty{2.01}{g/cm^3}$, the measured value.
\textbf{15.} Zinc blende: diamond is zinc blende with carbon on both kinds
of positions.
\textbf{16.} $a\sqrt3/4 = \qty{154.5}{pm}$; $\rho = 8 \times
12.01/(6.022 \times 10^{23} \times (3.567 \times 10^{-8})^3) =
\qty{3.52}{g/cm^3}$.
\textbf{17.} In diamond only half the \omterm{def:b1:crystals-metals:site}{tetrahedral sites} are filled, and by
atoms as large as those of the \omterm{def:b1:crystals-metals:lattice}{lattice}, which touch only along the bonds:
$C = 0.34$. In fluorite all the sites are filled and the small cations
fill the gaps between large anions: $C = 0.61$.
\textbf{18.} $40.08 + 2 \times 19.00 = \qty{78.08}{g/mol}$.
\textbf{19.} $4 \times 78.08/(6.022 \times 10^{23}) = \qty{5.186e-22}{g}$.
\textbf{20.} $(5.463 \times 10^{-8})^3 = \qty{1.630e-22}{cm^3}$.
\textbf{21.} One \ce{F-} missing per hundred cells removes $19.00/(100
\times 312.3)$ of the mass, 0.06\,\%: invisible at three figures (and a
\ce{Ca^{2+}} must go too, or the \omterm{def:b1:crystals-metals:crystal}{crystal} would not be neutral).
\textbf{22.} $\rho = 5.186 \times 10^{-22}/1.630 \times 10^{-22} =
\textbf{\qty{3.18}{g/cm^3}}$, exactly the measured density: the model of
the cell, built from one X-ray measurement of $a$, accounts for the mass
of the mineral.
\end{solution}
